% ===== Спільна преамбула навчального посібника (XeLaTeX) =====
\usepackage{fontspec}
\setmainfont{Libertinus Serif}
\setsansfont{Libertinus Sans}
\setmonofont{DejaVu Sans Mono}[Scale=0.80]
\usepackage{polyglossia}
\setdefaultlanguage{ukrainian}
\setotherlanguage{english}
\usepackage{microtype}

\usepackage[a5paper,inner=20mm,outer=13mm,top=17mm,bottom=18mm,headsep=4.5mm,footskip=9mm]{geometry}

\usepackage{amsmath,amssymb}
\usepackage{graphicx}
\usepackage{array}
\usepackage{tabularx}
\usepackage{longtable}
\usepackage{booktabs}
\usepackage{multirow}
\usepackage{enumitem}
\usepackage{framed}
\usepackage{tikz}
\usetikzlibrary{arrows.meta,positioning,shapes.geometric,fit,calc,backgrounds}
\usepackage{fancyhdr}
\usepackage{truncate}
\usepackage{titlesec}
\usepackage{emptypage}
\usepackage[hidelinks,unicode]{hyperref}
\usepackage{xurl}
\urlstyle{same}

\setlist{itemsep=1pt,parsep=1pt,topsep=3pt}
\setlength{\parindent}{1.25em}
\sloppy
\hbadness=10000
\vbadness=10000
\clubpenalty=10000
\widowpenalty=10000
\raggedbottom
\renewcommand{\arraystretch}{1.15}

% ---- Колонтитули та колонцифри (чергування парних/непарних сторінок) ----
\renewcommand{\chaptermark}[1]{\markboth{\thechapter.~#1}{}}
\renewcommand{\sectionmark}[1]{\markright{\thesection.~#1}}
\pagestyle{fancy}
\fancyhf{}
\newcommand{\chaptershort}[1]{\markboth{\thechapter.~#1}{}}
\fancyhead[LE]{\footnotesize\itshape\truncate{\headwidth}{\nouppercase{\leftmark}}}
\fancyhead[RO]{\footnotesize\itshape\hfill\truncate{\headwidth}{\nouppercase{\rightmark}}}
\fancyfoot[LE]{\footnotesize\bfseries\thepage}
\fancyfoot[RO]{\footnotesize\bfseries\thepage}
\renewcommand{\headrulewidth}{0.4pt}
\renewcommand{\footrulewidth}{0pt}
\fancypagestyle{plain}{%
  \fancyhf{}%
  \fancyfoot[LE]{\footnotesize\bfseries\thepage}%
  \fancyfoot[RO]{\footnotesize\bfseries\thepage}%
  \renewcommand{\headrulewidth}{0pt}}

% ---- Оформлення заголовків ----
\titleformat{\chapter}[display]
  {\normalfont\sffamily\bfseries\raggedright}
  {\normalsize\mdseries\MakeUppercase{Розділ~\thechapter}}
  {6pt}{\large}[\vspace{3pt}\titlerule]
\titlespacing*{\chapter}{0pt}{0pt}{16pt}
\titleformat{name=\chapter,numberless}[display]
  {\normalfont\sffamily\bfseries\raggedright}
  {}{0pt}{\large}[\vspace{3pt}\titlerule]
\titleformat{\section}{\normalfont\sffamily\bfseries\normalsize}{\thesection}{0.6em}{}
\titlespacing*{\section}{0pt}{11pt}{4pt}
\titleformat{\subsection}{\normalfont\sffamily\bfseries\small}{\thesubsection}{0.6em}{}
\titlespacing*{\subsection}{0pt}{8pt}{3pt}
\titleformat{\subsubsection}{\normalfont\sffamily\itshape\small}{\thesubsubsection}{0.6em}{}
\setcounter{secnumdepth}{3}
\setcounter{tocdepth}{1}

% ---- Авторські середовища та команди ----
\newenvironment{prnbox}{\par\smallskip\begin{leftbar}\small}{\end{leftbar}\smallskip\par}

\newcommand{\prnhead}{\section*{Відповідність програмним результатам навчання}%
  \addcontentsline{toc}{section}{Відповідність програмним результатам навчання}%
  \markright{Відповідність програмним результатам навчання}}

\newcommand{\selfcheck}{%
  \section*{Питання для самоперевірки}%
  \addcontentsline{toc}{section}{Питання для самоперевірки}%
  \markright{Питання для самоперевірки}}

\newcommand{\term}[1]{\textbf{#1}}
\newcommand{\figcap}[1]{\par\vspace{-2pt}\begin{center}\footnotesize #1\end{center}\vspace{2pt}}
% Рисунок разом із підписом (нерозривно)
\newenvironment{tikzfig}{\par\addvspace{6pt}\noindent\begin{minipage}{\linewidth}\centering}{\end{minipage}\par\addvspace{4pt}}
\newcommand{\tabcap}[1]{\par\vspace{2pt}\begin{center}\footnotesize #1\end{center}\vspace{2pt}}
%\newcommand{\tabhead}[1]{\par\vspace{4pt}\begin{center}\footnotesize #1\end{center}\vspace{-2pt}}
% Назва таблиці: мінімальний відступ до тіла таблиці й заборона розриву сторінки
% між назвою та таблицею (у т.ч. коли таблицю обгорнуто в center).
\usepackage{needspace}
\makeatletter
\newif\if@tabheadpending
\newif\if@tabheadincenter
\newcommand{\tabhead}[1]{%
  \par\addvspace{6pt}%
  \needspace{6\baselineskip}%
  {\centering\footnotesize\interlinepenalty\@M #1\par}%
  \nobreak\vskip 1.5pt\relax\nobreak
  \@nobreaktrue
  \global\@tabheadpendingtrue}
\AddToHook{env/center/begin}{%
  \if@tabheadpending
    \global\@tabheadpendingfalse\@tabheadincentertrue
    \topsep\z@\partopsep\z@
  \fi}
\AddToHook{env/center/end}{%
  \if@tabheadincenter\@topsepadd 6pt plus 2pt\relax\fi}
% таблиця одразу після \tabhead без center: без абзацного відступу, з інтервалом після
\newcommand{\@tabheadbefore}{\if@tabheadpending\ifvmode\noindent\fi\fi}
\newcommand{\@tabheadafter}{\if@tabheadpending
    \global\@tabheadpendingfalse\par\global\@nobreakfalse\addvspace{6pt plus 2pt}\fi}
\AddToHook{env/tabularx/before}{\@tabheadbefore}
\AddToHook{env/tabular/before}{\@tabheadbefore}
\AddToHook{env/tabularx/after}{\@tabheadafter}
\AddToHook{env/tabular/after}{\@tabheadafter}
\AddToHook{env/longtable/after}{\global\@tabheadpendingfalse}
\makeatother

\usepackage{ragged2e}
\newcolumntype{L}[1]{>{\RaggedRight\arraybackslash\hyphenpenalty=50\relax}p{#1}}
\newcolumntype{Y}{>{\RaggedRight\arraybackslash\hyphenpenalty=50\relax}X}
